% ch10 译文：原文 §6a 导出 dg 范畴、§6b 等变 dg 范畴（前）（main.tex 第 3291--3657 行）
\section{扭 Künneth 公式\label{sec:kunneth-formula}}
% [原文 p.41]
在引入等变层之后，还剩 Künneth 公式要建立，而它正是我们在定理 \ref{thm:twisted-tensor-P1-bundle} 的证明中实质上用到的工具。我们先建立拟凝聚层三角范畴的 dg 增强，它把凝聚层有界导出范畴与完美复形的 dg 增强作为 dg 全子范畴包含在内。

\subsection{导出 dg 范畴}

设 $X$ 是 $\field$ 上的概形，$\mathfrak{U}:=(U_i)_{i\in I}$ 为 $X$ 的开覆盖。对 $X$ 上任意阿贝尔群层 $\mathcal{H}$，回忆 $\mathcal{H}$ 关于 $\mathfrak{U}$ 的\emph{Čech 上链复形} $\check{C}^{\bullet}(\mathfrak{U},\mathcal{H})$ 由如下阿贝尔群族给出：
$$
\check{C}^p(\mathfrak{U},\mathcal{H})
:=
\prod_{i_0<\cdots<i_p}
\mathcal{H}\left(U_{i_0\cdots i_p}\right)
\quad\text{其中 }
p\in\mathbb{Z}_{\ge0},
$$
其中 $U_{i_0\cdots i_p}:=U_{i_0}\cap\cdots\cap U_{i_p}$；微分
$
\delta:
\check{C}^p(\mathfrak{U},\mathcal{H})
\rightarrow
\check{C}^{p+1}(\mathfrak{U},\mathcal{H})
$
定义为
$$
\left(\delta\alpha\right)_{i_0\cdots i_{p+1}}
:=
\sum_{j=0}^{p+1}
(-1)^j
\left.\alpha_{i_0\cdots\hat{i_j}\cdots i_{p+1}}\right|_{U_{i_0\cdots i_{p+1}}}
\in
\mathcal{H}(U_{i_0\cdots i_{p+1}})
$$
其中 $i_0<\cdots<i_{p+1}$，且记
$
\alpha
=
\left(
\alpha_{i_0\cdots i_p}
\in\mathcal{H}(U_{i_0\cdots i_p})
\right)_{i_0<\cdots<i_p}.
$

\begin{conceptbox}{Čech 复形}
给定开覆盖 $\mathfrak{U}=(U_i)_{i\in I}$ 与层 $\mathcal{H}$，Čech 复形在第 $p$ 次取所有 $p+1$ 重交上截面之积，微分 $\delta$ 交替删去一个下标并带符号 $(-1)^j$；当覆盖是仿射开覆盖而层拟凝聚时，它的上同调就是层上同调 $H^{\bullet}(X,\mathcal{H})$。这里取乘积而不取直和是必需的：一个整体截面在每个 $U_i$ 上都有分量，覆盖无限时它并不属于直和。本文用它把每对对象之间的 $\mathcal{H}om$ 复形拼装成 dg 态射复形，从而造出 $\Coh(X;\mathfrak{U})_{\dg}$。
\par\smallskip\noindent{\footnotesize 参考：Hartshorne, \emph{Algebraic Geometry}, GTM 52, Springer, 1977, 第三章；本文附录 A。}
\end{conceptbox}

\begin{yizhu}[Čech 微分里的符号与删除]
$\delta$ 的第 $j$ 项是「删去下标 $i_j$、再限制到更小的交」，符号 $(-1)^j$ 依删去位置交替，它保证 $\delta^2=0$。另一个易错处是 $\check{C}^p$ 取乘积 $\prod$ 而非直和：直和版本在无限覆盖下连「整体截面的像」都装不下，$\check{H}^0$ 会算错。
\end{yizhu}

\begin{dfn}\label{dfn:dg-enhancement-Coh}
设 $X$ 是存在仿射开覆盖 $\mathfrak{U}:=(U_i)_{i\in I}$ 的概形。$X$ 关于 $\mathfrak{U}$ 的\emph{凝聚导出 dg 范畴}是一个 dg 范畴
$$
\Coh(X;\mathfrak{U})_{\dg}
$$
其定义如下：
\begin{itemize}
\item
\emph{对象}由凝聚 $\mathcal{O}_X$-模的有界上链复形
$$
\mathcal{F}^{\bullet}
=
\left(
\bigoplus_{\objectdegree\in\mathbb{Z}}
\mathcal{F}^{\objectdegree},
d_{\mathcal{F}}:=\left(d_{\mathcal{F}}^{\objectdegree}:\mathcal{F}^{\objectdegree}\rightarrow \mathcal{F}^{\objectdegree+1}\right)_{\objectdegree\in\mathbb{Z}}\right)
$$
给出，其中 $\mathcal{F}^{\objectdegree}\in\Coh(X)$。也就是说，微分 $d_{\mathcal{F}}=(d_{\mathcal{F}}^{\objectdegree})_{\objectdegree\in\mathbb{Z}}$ 满足 $d_{\mathcal{F}}^{\objectdegree}\in \Hom_{\Coh(X)}(\mathcal{F}^{\objectdegree},\mathcal{F}^{\objectdegree+1})$，且对每个 $\objectdegree\in\mathbb{Z}$ 都有 $d_{\mathcal{F}}^{\objectdegree}\comp d_{\mathcal{F}}^{\objectdegree+1}=0$。

\item
给定两个对象 $\mathcal{F}$ 与 $\mathcal{F}'$，先定义凝聚 $\mathcal{O}_X$-模的有界上链复形
$$
\mathcal{H}om_{\mathcal{O}_X}^q\left(\mathcal{F},\mathcal{F}'\right)
:=
\bigoplus_{\objectdegree\in\mathbb{Z}}
\mathcal{H}om_{\mathcal{O}_X}\left(\mathcal{F}^{\objectdegree},\mathcal{F}'^{\objectdegree+q}\right)
\quad\text{其中 }q\in\mathbb{Z},
$$
其微分
$
d_{\mathcal{H}om}:
\mathcal{H}om_{\mathcal{O}_X}^q(\mathcal{F},\mathcal{F}')
\rightarrow
\mathcal{H}om_{\mathcal{O}_X}^{q+1}(\mathcal{F},\mathcal{F}')
$
由
$$
d_{\mathcal{H}om}f:=d_{\mathcal{F}}\comp f - (-1)^q f\comp d_{\mathcal{F}'}
$$
定义。

$\mathcal{F}$ 到 $\mathcal{F}'$ 的\emph{dg-态射空间}是 $\field$-模的上链复形，其第 $k$ 次分量由
\begin{align*}
\hom_{\Coh(X;\mathfrak{U})_{\dg}}^k
\left(
\mathcal{F},\mathcal{F}'
\right)
:&=
\bigoplus_{p+q=k}
\check{C}^p\big(\mathfrak{U},\mathcal{H}om_{\mathcal{O}_X}^q
(\mathcal{F},\mathcal{F}')
\big),
\\[2mm]
&=
\bigoplus_{p+q=k}
\prod_{i_0<\cdots<i_p}\bigoplus_{\objectdegree\in\mathbb{Z}}
\Hom_{\Coh(U_{i_0\cdots i_p})}\left(\mathcal{F}^{\objectdegree}|_{U_{i_0\cdots i_p}},\mathcal{F}'^{\objectdegree+q}|_{U_{i_0\cdots i_p}}\right)
\end{align*}
给出，而\emph{微分}取全微分
$$
D
:=
\delta + (-1)^p d_{\mathcal{H}om}.
$$

\item
\emph{复合}定义为
\begin{align*}
(f\comp f')_{i_0\cdots i_{p+p'}}^{\objectdegree}
:&=
% [原文 p.42]
(-1)^{qp'}
\left.
f_{i_0\cdots i_p}^{\objectdegree}
\right|_{U_{i_0\cdots i_{p+p'}}}
\comp
\left.
{f'}_{i_p\cdots i_{p+p'}}^{\objectdegree+q}
\right|_{U_{i_0\cdots i_{p+p'}}}
\\[3mm]
&\in
\Hom_{\Coh(U_{i_0\cdots i_{p+p'}})}\left(\mathcal{F}^{\objectdegree}|_{U_{i_0\cdots i_{p+p'}}},\mathcal{F}''^{\objectdegree+q+q'}|_{U_{i_0\cdots i_{p+p'}}}\right)
\end{align*}
其中
\begin{align*}
f
&:=
\left(
f_{i_0\cdots i_p}^{\objectdegree}
\in
\Hom_{\Coh(U_{i_0\cdots i_p})}\left(\mathcal{F}^{\objectdegree}|_{U_{i_0\cdots i_p}},\mathcal{F}'^{\objectdegree+q}|_{U_{i_0\cdots i_p}}\right)
\right)_{
\substack{
i\in\mathbb{Z}
\\
i_0<\cdots<i_p
}
}
\\[2mm]
&\in
\check{C}^p\big(
\mathfrak{U},\mathcal{H}om_{\mathcal{O}_X}^q(\mathcal{F},\mathcal{F}')
\big),
\\[2mm]
f'
&:=
\left(
{f'}_{i_0\cdots i_{p'}}^{\objectdegree}
\in
\Hom_{\Coh(U_{i_0\cdots i_{p'}})}\left(\mathcal{F}'^{\objectdegree}|_{U_{i_0\cdots i_{p'}}},\mathcal{F}''^{\objectdegree+q'}|_{U_{i_0\cdots i_{p'}}}\right)
\right)_{
\substack{
i\in\mathbb{Z}
\\
i_0<\cdots<i_p
}
}
\\[2mm]
&\in
\check{C}^{p'}\big(
\mathfrak{U},\mathcal{H}om_{\mathcal{O}_X}^{q'}(\mathcal{F}',\mathcal{F}'')
\big).
\end{align*}

\item
$\mathcal{F}$ 的\emph{单位}
$
\id_{\mathcal{F}}
:=
\left(\left(\id_{\mathcal{F}}\right)_{i_0\cdots i_p}^{\objectdegree}\right)_{
\substack{
i\in\mathbb{Z}
\\
i_0<\cdots<i_p
}
}
\in
\check{C}^p\big(
\mathfrak{U},\mathcal{H}om_{\mathcal{O}_X}^q(\mathcal{F},\mathcal{F})
\big)
$
由
$$
\left(\id_{\mathcal{F}}\right)_{i_0}^{\objectdegree}
:=
\id_{\mathcal{F}^{\objectdegree}|_{U_{i_0}}}
\in
\Hom_{\Coh(U_{i})}\left(\mathcal{F}^{\objectdegree}|_{U_{i_0}},\mathcal{F}^{\objectdegree}|_{U_{i_0}}\right)
\quad\text{对每个 }
i_0\in I,\
\objectdegree\in\mathbb{Z}
$$
给出，其中 $p=0$、$q=0$，其余情形取零。
\end{itemize}

这些运算满足 $D^2=0$、结合律，以及 Leibniz 法则
$$
D(f\comp f')
=
Df\comp f'+(-1)^{|f|}f\comp Df'
$$
其中 $f$ 与 $f'$ 为齐次态射。
\end{dfn}

\begin{yizhu}[指标 $i\in\mathbb{Z}$ 是原文笔误]
复合与单位的公式里，求积下标集合写作 $i\in\mathbb{Z}$，按定义应为 $i_0<\cdots<i_p$ 遍历覆盖指标集 $I$；同一处既写 $i$ 又写 $i_0<\cdots<i_p$，显系排印之误。照抄公式时不必顺手改成 $i\in I$（否则与原文对照时找不到出处），阅读时按 $I$ 理解即可。
\end{yizhu}

按惯例，把完美复形（局部\emph{同构}于向量丛有界复形者）构成的 dg 全子范畴记作 $\Perf (X;\mathfrak{U})_{dg}$。

\begin{yizhu}[$\Perf$ 的约定与通常不同]
这里的 $\Perf$ 指局部完美（locally perfect），而非要求全局向量丛复形代表的严格完美（strictly perfect）；原文用注释说明了这一点，但这层区别在译文里只能靠这条注保留。$X$ 满足 (GSP) 条件时两种定义在 $D\QCoh(X)$ 中一致，本文取局部版本，正是为了绕开对该条件的额外假设。
\end{yizhu}

\begin{prp}\label{prp:verifying-dg}
    $\Coh (X;\mathfrak{U})_{dg}$ 是 dg 范畴。
\end{prp}
\begin{proof}
    只需验证复合 $\comp$ 的结合律。这可以直接验证，但要十分小心符号。
\end{proof}

\begin{prp}\label{prp:dg-enhancement}
$\Coh(X,\mathfrak{U})_{\dg}$ 是 $D^b(\Coh(X))$ 的 dg 增强。
\end{prp}
\begin{proof}
    回忆松弛 $\mathcal{O}_X$-模 $\mathcal{F}$ 满足 $\check{H}^p (\mathfrak{U},\mathcal{F})=0$（$p>0$）（\cite[Chapter 3, Proposition 4.3]{Hartshorne77}），且内射 $\mathcal{O}_X$-模都是松弛的（\cite[Chapter 3, Lemma 2.4]{Hartshorne77}）。由于仿射簇上的整体截面是正合函子，Čech 复形 $\check{C}^{\bullet} (\mathfrak{U};-)\colon\Coh (\mathcal{O}_X)\to\Ch (\Bbbk)$ 作为阿贝尔范畴之间的函子是正合的。\par
    给定层复形 $\mathcal{F}^{\bullet}\in\Ch (\mathcal{O}_X)$，设 $\mathcal{F}^{\bullet}\to\mathcal{I}^{\bullet,\bullet}$ 为其内射消解构成的双复形，则由正合性可知 $\check{C}^{\bullet} (\mathfrak{U};\mathcal{F}^{\bullet})\to\check{C}^{\bullet} (\mathfrak{U};\mathcal{I}^{\bullet,\bullet})$ 是链复形的拟同构。因此可以把 $\mathcal{F}^{\bullet}$ 换成内射消解的全复形 $\Tot^{\bullet} (\mathcal{I}^{\bullet,\bullet})$。
% [原文 p.43]
    此时 Čech 复形 $\check{C}^{\bullet}\left(\Tot^{\bullet} (\mathcal{I}^{\bullet,\bullet})\right)$ 可以写成一个双复形。利用上述内射 $\mathcal{O}_X$-模的性质可知，对任意 $(p,q)$，$\displaystyle\check{H}^{\bullet} (\mathcal{I}^{p,q})$ 都集中在 $0$ 次。双复形相联的谱序列 \cite[\href{https://stacks.math.columbia.edu/tag/012X}{Section 012X}]{stacks-project} 给出谱序列 $E^{p,q}_d$，使得 $E_2^{p,q}=\displaystyle\check{H}^q\left(\mathfrak{U};H^p(\Tot^{\bullet}(\mathcal{I}^{\bullet,\bullet}))\right)$ 集中在 $q=0$ 处，故谱序列在第一页退化，从而有同构\[
    H^p\left(\Tot^{\bullet}(\check{C}^{\bullet} (\mathfrak{U};\Tot^{\bullet} (\mathcal{I}^{\bullet,\bullet})))\right)\cong H^p (\Tot^{\bullet} (\mathcal{I}^{\bullet,\bullet}))\cong H^p (\mathcal{F}^{\bullet}).
    \]
    现在对任意对象 $\mathcal{F}^{\bullet}$ 与 $\mathcal{G}^{\bullet}$，有 $\field$-模的同构\[
    H^0\hom^{\bullet}_{\Coh (X)} (\mathcal{F}^{\bullet},\mathcal{G}^{\bullet})_{dg}\simeq\bigoplus_{p\in\mathbb{Z}}\hom_{\Coh (X)} (\mathcal{F}^p,\mathcal{G}^p).
    \]再结合 $H\Mod (\Coh(X)_{dg}))$ 中的平移 $[1]$ 保持 $\Coh (X;\mathfrak{U})_{dg}$ 这一事实，即完成证明。
\end{proof}

\begin{conceptbox}{dg 增强}
三角范畴只记得态射空间的上同调，丢掉了把 $\Hom$ 提升为复形这一层信息；\emph{dg 增强}就是找一个 dg 范畴，其同伦范畴与原三角范畴等价，且 $\hom^{\bullet}$ 复形实现给定的导出 $\Hom$。有了它，才能谈 dg 函子、$A_\infty$ 函子以及扭 Künneth 公式这类需要高阶结构的构造。本文的 $\Coh(X;\mathfrak{U})_{\dg}$ 与 $\Coh_G(X;\mathfrak{U})_{\dg}$ 分别是 $\Coh$ 与 $\Coh_G$ 的 dg 化；命题 \ref{prp:dg-enhancement} 说明前者确实是 $D^b(\Coh(X))$ 的增强。
\par\smallskip\noindent{\footnotesize 参考：nLab，\emph{dg-category}；Keller, \emph{On differential graded categories}, arXiv:math/0601185。}
\end{conceptbox}

\begin{yizhu}[谱序列为什么退化为同构]
关键在 $E_2^{p,q}=\check{H}^q\left(\mathfrak{U};H^p(\Tot^{\bullet}(\mathcal{I}^{\bullet,\bullet}))\right)$ 只支撑在 $q=0$：内射模松弛，而松弛层在开覆盖上的 Čech 上同调在正次数消失。于是所有 $q\neq0$ 的项为零，谱序列在第一页之后不再变化，$H^p$ 的总复形上同调便直接同构于 $H^p$ 的全复形上同调——这就是「用内射消解计算导出函子」在 Čech 情形下的落地方式。
\end{yizhu}

\subsection[等变 dg 范畴]{$G$-等变导出 dg 范畴}

设 $G$ 是仿射群概形，$\mathfrak U:=(U_i)_{i\in I}$ 是 $X$ 的仿射开覆盖，且每个 $U_i$ 都是 $G$-不变的。对 $G$-等变 $\mathcal O_X$-模 $(\mathcal H,\Theta)$，它关于 $\mathfrak U$ 的\emph{$G$-不变 Čech 上链复形}由如下阿贝尔群族给出：
$$
\check C_G^p\big(\mathfrak U,(\mathcal H,\Theta)\big)
:=
\prod_{i_0<\cdots<i_p}
\mathcal H(U_{i_0\cdots i_p})^G
\quad\text{其中 }p\in\mathbb{Z}_{\ge0},
$$
其中 $U_{i_0\cdots i_p}:= U_{i_0}\cap\cdots\cap U_{i_p}$，微分
$
\delta:
\check{C}_G^p\big(\mathfrak{U},(\mathcal{H},\Theta)\big)
\rightarrow
\check{C}_G^{p+1}\big(\mathfrak{U},(\mathcal{H},\Theta)\big)
$
定义为
$$
(\delta s)_{i_0\cdots i_{p+1}}
:=
\sum_{j=0}^{p+1}(-1)^j
\left.
s_{i_0\cdots\widehat{i_j}\cdots i_{p+1}}
\right|_{U_{i_0\cdots i_{p+1}}}
\in
\mathcal{H}(U_{i_0\cdots i_{p+1}})^G
$$
其中 $i_0<\cdots<i_{p+1}$，且记
$
s=\left(s_{i_0\cdots i_p}\in\mathcal{H}(U_{i_0\cdots i_p})^G\right)_{i_0<\cdots<i_p}$。
